% broader impacts: drafted 2026-07-19 via the style-engine protocol for the
% NeurIPS submission restructure. Unnumbered section placed after the
% Conclusion, before References. No new numbers: every figure restates a
% value already established in the paper.
Prediction markets are zero-sum. Every dollar the strategy wins is lost by
other bettors in the same book, mostly retail traders who price the open as
if the futures gap carries in full. Publishing the mispricing is the fairer
move: once the $k \approx 0.8$ pass-through is public, the crowd can
reprice it, and the competition scan in Section~\ref{sec:discussion}
suggests informed behavior does get copied. The harm is also bounded by the
market itself. The stake ceiling is 1\% of a \$30--48k daily book, near
\$1{,}000 per day, so no scaled version of this system can extract more
than the thin volume the market brings, and that capacity falls one-for-one
as volume falls.

The gambling adjacency is real and we do not dress it up. These positions
are bets that resolve within hours, and the comment threads quoted in
Section~\ref{sec:introduction} record a bettor losing half of his money on
one quirk day. Nothing in this paper improves the odds of an unmodeled
bettor; it documents that the crowd's default belief loses to a measured
one. The responsible use is the system's own rule: bet only with a modeled
edge, size by a rule that caps ruin, and stand down on coin-toss days. A
reader without a model has no edge to trade, and Section~\ref{sec:results}
shows that even ours is promising but statistically unconfirmed.
